\paragraph{The projection formula and its missing inputs.}
The current literature supplies another precise test for a proposed
orbit-knot computation. Theorem 4.1 of
\citet{GauniyalTurchin2026Root}, published online on
25 September 2026, expresses the Haefliger integer using both
double-point linking and invariants of the projected immersion.
The latter include a rationally normalized Smale invariant.
Consequently the double circle of C.3.17 alone is insufficient.
An application here would require an actual compressed projection
and its tangent-framing comparison. We opened the published
formula and its definitions, rather than importing the linking
term in isolation.

Likewise, \citet[Theorem 2.5]{Sakai2015Root} computes the
\emph{change} under a specified crossing change from oriented
double-preimage linking. To obtain an absolute value by that
route one must also identify the comparison embedding and its
class. Neither formula identifies the present trace with a
standard reference embedding. The characteristic-number cap
retains those geometric comparisons as explicit requirements.
Searches on 15 November 2026 for the six-sphere family,
its degeneration monodromy, and these knot formulas did not
locate a numerical determination for this particular family.
This states the search outcome and claims no priority.
